

\chapter{答案-不等式}

\section{不等式}


\subsection{朱韬数学-不等式-1}
\label{朱韬数学-不等式-1}

\begin{center}
\includegraphics[width=0.4\textwidth]{朱韬数学-不等式-1.png}
\end{center}


\subsection{2024北京东城期末}
\label{2024北京东城期末}

\textbf{第一步：解不等式组}  
由第二个不等式：
\[
2x - 3 < 1 \implies 2x < 4 \implies x < 2.
\]
结合第一个不等式 \(x \ge a\)，得到解集为
\[
a \le x < 2.
\]

\textbf{第二步：确定整数解}  
因为解集是左闭右开的区间 \([a, 2)\)，其中右端点 \(2\) 不包含在内，所以可能的整数解只能从小于 \(2\) 的整数中选取。  
小于 \(2\) 的整数有 \(\dots, -2, -1, 0, 1\)。

若要使区间内恰好有 \(3\) 个整数，则这 \(3\) 个整数必须是连续的。由于最大整数不能超过 \(1\)（因为 \(x<2\)），且区间是左闭的，所以三个整数只可能是：
\[
-1,\;0,\;1.
\]

\textbf{第三步：由端点确定 \(a\) 的范围}  
要保证 \(-1, 0, 1\) 都在解集内，且 \(-2\) 不在解集内（否则会有四个整数解），同时 \(2\) 已排除。

- 最小的整数解是 \(-1\)，所以必须有 \(a \le -1\)（才能包含 \(-1\)）。
- 要排除 \(-2\)，必须有 \(a > -2\)（因为若 \(a \le -2\)，则 \(-2\) 也会满足 \(x \ge a\)，从而进入解集，整数解将至少为 \(-2, -1, 0, 1\) 共四个）。

因此：
\[
-2 < a \le -1.
\]


\subsection{2025北京二中期末}
\label{2025北京二中期末}

第（1）问

化简不等式：
\[
|x| + 3 < 2 \implies |x| < -1.
\]
由于绝对值恒为非负，不可能小于负数，因此该不等式无解，其解集为空集。而“邻域”定义为开区间 \((x_0-d, x_0+d)\)，空集无法表示成这种形式，所以不存在这样的邻域。

\[
\boxed{\text{不存在}}
\]

第（2）问

首先根据邻域定义写出两个不等式的解集。

不等式 \(A\) 的解集构成“\(a+1\) 的 \(2\) 邻域”，即
\[
A:\quad (a+1-2,\; a+1+2) = (a-1,\; a+3).
\]

不等式 \(B\) 的解集构成“\(2\) 的 \(1\) 邻域”，即
\[
B:\quad (2-1,\; 2+1) = (1,\; 3).
\]

不等式组由 \(A\) 和 \(B\) 组成，其解集为两者的交集：
\[
A \cap B = (a-1,\; a+3) \cap (1,\; 3).
\]
交集的一般形式为
\[
\bigl(\max(a-1,1),\; \min(a+3,3)\bigr).
\]

又已知该交集恰好构成“\(\left(\dfrac{a}{2}+2\right)\) 的 \(\left(\dfrac{a}{2}+1\right)\) 邻域”，即
\[
\left(\left(\frac{a}{2}+2\right)-\left(\frac{a}{2}+1\right),\;
\left(\frac{a}{2}+2\right)+\left(\frac{a}{2}+1\right)\right)
= (1,\; a+3).
\]
注意邻域要求半径 \(d = \dfrac{a}{2}+1 > 0\)，故
\[
a > -2. \tag{1}
\]

因此必须有
\[
\max(a-1,1) = 1 \quad \text{且} \quad \min(a+3,3) = a+3.
\]

由 \(\max(a-1,1)=1\) 得
\[
a-1 \le 1 \implies a \le 2. \tag{2}
\]

由 \(\min(a+3,3)=a+3\) 得
\[
a+3 \le 3 \implies a \le 0. \tag{3}
\]

结合 (1)、(2)、(3)，取交集：
\[
-2 < a \le 0.
\]

同时还需保证交集非空且右端点大于左端点，即 \(1 < a+3\)，这等价于 \(a > -2\)，已包含在 (1) 中。

因此 \(a\) 的取值范围为
\[
\boxed{-2 < a \le 0}.
\]

{答案}

\begin{enumerate}
    \item 不存在；
    \item \(-2 < a \le 0\)。
\end{enumerate}


\subsection{2024北京房山期末}
\label{2024北京房山期末}

设甲、乙、丙所写数字分别为 \(a, b, c\)，均为正整数。

将条件转化为不等式：

\begin{enumerate}
    \item 乙的一半大于甲：\(\dfrac{b}{2} > a\)，即 \(b > 2a\)。
    \item 丙不大于甲：\(c \le a\)。
    \item 丙的 3 倍大于乙：\(3c > b\)，即 \(b < 3c\)。
    \item \(a, b, c \in \mathbb{N}^*\)（正整数）。
\end{enumerate}

综合得：
\[
2a < b < 3c \quad \text{且} \quad c \le a.
\]

由 \(b > 2a\) 和 \(b < 3c\) 得：
\[
2a < 3c.
\]
又因为 \(c \le a\)，所以 \(3c \le 3a\)，但我们需要的是 \(2a < 3c\)，这给出 \(a\) 和 \(c\) 之间的大致关系。

由于 \(a, c\) 为正整数，且 \(c \le a\)，尝试从最小的 \(a\) 开始枚举。

\textbf{若 \(a=1\)}：则 \(c \le 1\)，故 \(c=1\)。  
此时 \(b > 2\) 且 \(b < 3\)，即 \(2 < b < 3\)，无正整数解。

\textbf{若 \(a=2\)}：则 \(c \le 2\)，\(c\) 可取 1 或 2。  
- 当 \(c=1\) 时，\(b > 4\) 且 \(b < 3\)，矛盾。  
- 当 \(c=2\) 时，\(b > 4\) 且 \(b < 6\)，故 \(b\) 只能取 5。  
此时 \((a,b,c)=(2,5,2)\)，和为 \(2+5+2=9\)。

\textbf{检查是否有更小的和}：\(a=1\) 无解，最小和至少要从 \(a=2\) 开始，我们得到一组可行解和为 9。但还需验证是否存在和为更小的情况（比如和为 7 或 4），由于 \(a\ge1\)，若和为 7，则可能组合如 \((1,?,?)\) 或 \((2,?,?)\) 但 \(a=1\) 无解，\(a=2\) 时最小可能和为 \(2+5+2=9\)，所以 7 不可能。因此最小和为 9。

故答案为 C。


\subsection{不等式组含参问题-基础2}

\label{不等式组含参问题-基础2}  % 为答案设置标签，记录当前页码

不等式 \((1+2a)x > 1\) 的解集为 \(x < \dfrac{1}{1+2a}\)。

这表明在解不等式时，不等号方向发生了改变。不等式两边同除以一个负数时，不等号方向才会改变，因此必须满足：
\[
1+2a < 0.
\]
解得：
\[
2a < -1 \quad\Longrightarrow\quad a < -\dfrac{1}{2}.
\]

因此 \(a\) 的取值范围为 \(a < -\dfrac{1}{2}\)，对应选项 \textbf{A}。

\vspace{2mm}
\noindent 答案：\textbf{A}







\section{不等式（组）与方程（组）的综合应用}

\subsection{2025北京二中期末}
\label{2025北京二中期末}
\subsection*{(1)}
由 \(x+y=3\)，当 \(x=0\) 时，\(y=3\)，故
\[
p=3.
\]

\subsection*{(2)}
分别联立各方程与 \(x+y=3\)，并检查解是否满足 \(-7<x<8\)。

\begin{itemize}
    \item ① \(3x-y=1\)，代入 \(y=3-x\)：
    \[
    3x-(3-x)=1 \Rightarrow 4x=4 \Rightarrow x=1,
    \]
    满足 \(-7<1<8\)，有解。

    \item ② \(2x+y=-5\)，代入 \(y=3-x\)：
    \[
    2x+3-x=-5 \Rightarrow x=-8,
    \]
    不满足范围，无解。

    \item ③ \(x+2y=-4\)，代入 \(y=3-x\)：
    \[
    x+2(3-x)=-4 \Rightarrow -x+6=-4 \Rightarrow x=10,
    \]
    不满足范围，无解。
\end{itemize}

所以填：
\[
\boxed{①}
\]

\subsection*{(3)}
方程 \(cx+dy=1\) 的部分解为 \((-7,-2)\)，\((2,n)\)，\((8,n+6)\)。

由斜率相等：
\[
\frac{n-(-2)}{2-(-7)} = \frac{(n+6)-n}{8-2}
\]
即
\[
\frac{n+2}{9} = \frac{6}{6} = 1,
\]
解得
\[
n=7.
\]

于是直线过 \((-7,-2)\) 和 \((2,7)\)，代入 \(cx+dy=1\) 得：
\[
-7c-2d=1,\qquad 2c+7d=1.
\]
解得
\[
c=-\frac15,\quad d=\frac15.
\]
所以该方程为
\[
-\frac15 x + \frac15 y = 1 \implies y=x+5.
\]

与 \(x+y=3\) 联立：
\[
\begin{cases}
x+y=3,\\
y=x+5
\end{cases}
\]
解得
\[
x=-1,\quad y=4.
\]

故方程组解为
\[
\boxed{x=-1,\; y=4}.
\]



\subsection{2025北京门头沟期末}
\label{2025北京门头沟期末}
\subsection*{（1）}

设甲种T恤每件进价为 \(x\) 元，乙种T恤每件进价为 \(y\) 元。

由题意得：
\[
\begin{cases}
2x + 3y = 310, \\
x + 2y = 190.
\end{cases}
\]

解方程组：

由第二个方程得 \(x = 190 - 2y\)，代入第一个方程：
\[
2(190 - 2y) + 3y = 310 \implies 380 - 4y + 3y = 310 \implies -y = -70 \implies y = 70.
\]
则
\[
x = 190 - 2 \times 70 = 190 - 140 = 50.
\]

所以甲种T恤每件进价 \(50\) 元，乙种T恤每件进价 \(70\) 元。

\[
\boxed{\text{甲：50元/件，乙：70元/件}}
\]

\subsection*{（2）}

设购进甲种T恤 \(m\) 件，则乙种T恤购进 \((100-m)\) 件。

根据题意，总费用不超过6540元：
\[
50m + 70(100-m) \le 6540.
\]
化简：
\[
50m + 7000 - 70m \le 6540 \implies -20m \le -460 \implies m \ge 23.
\]

又因为甲种数量的三倍不超过乙种数量：
\[
3m \le 100 - m \implies 4m \le 100 \implies m \le 25.
\]

同时 \(m\) 为正整数，且 \(0 \le m \le 100\)，结合上述得：
\[
23 \le m \le 25.
\]

因此 \(m\) 可取 \(23, 24, 25\)。

对应方案：

\begin{itemize}
    \item 当 \(m=23\) 时，甲种23件，乙种 \(100-23=77\) 件；
    \item 当 \(m=24\) 时，甲种24件，乙种 \(100-24=76\) 件；
    \item 当 \(m=25\) 时，甲种25件，乙种 \(100-25=75\) 件。
\end{itemize}

所以共有三种购买方案：

\[
\boxed{
\begin{array}{c|c|c}
\text{方案} & \text{甲种（件）} & \text{乙种（件）} \\
\hline
1 & 23 & 77 \\
2 & 24 & 76 \\
3 & 25 & 75
\end{array}
}
\]

\section{新定义与阅读理解型不等式}

\subsection{2025北京二中期末2}
\label{2025北京二中期末2}

对于实数 \(a, b\)，定义 \(\max\{a, b\}\)：当 \(a \geq b\) 时取 \(a\)，当 \(a < b\) 时取 \(b\)。

已知 \(\max\{2x + 1,\; 3x - 1\} = 3\)，求 \(x\) 的值。

\section*{解析}

根据定义，\(\max\{2x+1,\;3x-1\}\) 等于两者中较大的数。要使其等于 \(3\)，分两种情况讨论：

\subsection*{情况一：\(2x+1 \geq 3x-1\)}

此时最大值取 \(2x+1\)，故有
\[
2x+1 = 3 \implies 2x = 2 \implies x = 1.
\]
检验条件：\(2x+1 \geq 3x-1\) 即
\[
2(1)+1 \geq 3(1)-1 \implies 3 \geq 2,
\]
成立，所以 \(x=1\) 是可行解。

\subsection*{情况二：\(2x+1 < 3x-1\)}

此时最大值取 \(3x-1\)，故有
\[
3x-1 = 3 \implies 3x = 4 \implies x = \frac{4}{3}.
\]
检验条件：\(2x+1 < 3x-1\) 即
\[
2\cdot\frac{4}{3}+1 < 3\cdot\frac{4}{3}-1 \implies \frac{8}{3}+1 < 4-1 \implies \frac{11}{3} < 3,
\]
显然不成立（因为 \(\frac{11}{3} \approx 3.67 > 3\)），故 \(x=\frac{4}{3}\) 不符合前提，舍去。

答案：\[
\boxed{x = 1}
\]



\subsection{新定义-集合思想}

\label{新定义-集合思想}  % 为答案设置标签，记录当前页码

(1) 解方程组 $\begin{cases} 2x - y = 5 \\ 3x + 4y = 2 \end{cases}$ 得 $\begin{cases} x = 2 \\ y = -1 \end{cases}$，故 $C = \{2, -1\}$。  
不等式 $x + 1 \geq 0$ 的解集 $D = \{x \mid x \geq -1\}$。  
因为 $2 \geq -1$ 且 $-1 \geq -1$，所以 $C$ 中的所有元素都在 $D$ 中，故 $C$ 能被 $D$ 包含。

(2) 解方程组 $\begin{cases} 2x + 3y - 5a = -1 \\ x - 2y + a = 3 \end{cases}$ 得 $\begin{cases} x = a + 1 \\ y = a - 1 \end{cases}$，故 $E = \{a+1,\ a-1\}$。  
解不等式组 $\begin{cases} 3(x-2) \geq x-4 \\ 2x+1 \geq x-1 \end{cases}$：  
第一个得 $x \geq 1$，第二个得 $x \geq -2$，取交集得 $F = \{x \mid x \geq 1\}$。  
$E$ 不能被 $F$ 包含，即 $a+1 < 1$ 或 $a-1 < 1$（只要有一个数小于1即可，因为 $F$ 要求 $\geq 1$）。  
注意 $a+1 > a-1$，所以只需考虑较小者 $a-1 < 1$ 即可保证至少一个数不在 $F$ 中。解得 $a < 2$。  
或者也可以按原答案思路：$a+1 < 1$ 得 $a < 0$；$a-1 < 1$ 得 $a < 2$。综合取并集得 $a < 2$。  
因此 $a$ 的取值范围是 $a < 2$。

\subsection{二中教育集团期末第26题}

\label{二中教育集团期末第26题}  % 为答案设置标签，记录当前页码

(1) 解方程组 $\begin{cases} x + y = 4 \\ x = 1.5 \end{cases}$ 得 $\begin{cases} x = 1.5 \\ y = 2.5 \end{cases}$。  
在数轴上实数 $1.5$ 所对的点为 $A$，$2.5$ 所对的点为 $B$。线段 $AB$ 上存在 $1$ 个整数，故是 $1$ 系方程。

(2) 解方程组 $\begin{cases} x + y = 2.5 + 2m \\ x - 2y = 2.5 - 6m \end{cases}$ 得 $\begin{cases} x = 2.5 \\ y = 2m \end{cases}$。  
在数轴上实数 $2.5$ 所对的点为 $A$，$2m$ 所对的点为 $B$。  
该方程组是 $3$ 系方程，即在线段 $AB$ 上存在 $3$ 个整数。  
- 若整数为 $3,4,5$，则 $5 \le 2m < 6$，解得 $2.5 \le m < 3$。  
- 若整数为 $0,1,2$，则 $-1 < 2m \le 0$，解得 $-0.5 < m \le 0$。  
综上，$m$ 的取值范围为 $-0.5 < m \le 0$ 或 $2.5 \le m < 3$。

(3) 解方程组 $\begin{cases} x - y = 2n - 1.6 \\ x = n + 0.4 \end{cases}$ 得 $\begin{cases} x = n + 0.4 \\ y = 2 - n \end{cases}$。  
在数轴上实数 $n+0.4$ 所对的点为 $A$，$2-n$ 所对的点为 $B$。  
该方程组是 $2$ 系方程，即线段 $AB$ 上存在 $2$ 个整数，因此 $1 \le | (n+0.4) - (2-n) | < 3$，且区间端点的整数个数满足条件。  
化简得 $1 \le |2n - 1.6| < 3$，即 $1 \le |2n-1.6| < 3$。  
解得 $n$ 的取值范围为 $(-0.7, -0.3] \cup [1.3, 2.3)$（具体计算过程略）。


\section{不等式与最值综合}
\subsection{不等式与最值综合-题目1}
\label{不等式与最值综合-题目1}

分析：解不等式组得 \(-3 \leq x \leq 4\)，根据两点间的距离公式知当 \(-1 \leq x \leq 1\) 时，\(S = |x-1|+|x+1|\) 取得最小值；当 \(x=4\) 时，\(S\) 取最大值，继而可得答案。

详解：
解不等式 \(5(x+1) \geq 3x-1\)，得：
\[
x \geq -3,
\]
解不等式 \(\frac{1}{2}x-1 \leq 7-\frac{3}{2}x\)，得：
\[
x \leq 4,
\]
则
\[
-3 \leq x \leq 4.
\]

当 \(-1 \leq x \leq 1\) 时，\(S = |x-1|+|x+1|\) 取得最小值，最小值 \(n=2\)；  
当 \(x=4\) 时，\(S = |x-1|+|x+1|\) 取得最大值，最大值 \(m=8\)。  
因此
\[
mn = 2 \times 8 = 16.
\]

故答案为：\(\boxed{16}\)。


\subsection{不等式与坐标系-题目1}
\label{不等式与坐标系-题目1}

\begin{center}
\includegraphics[width=0.8\textwidth]{不等式与坐标系-题目1.png}
\end{center}






